\documentclass[xelatex,ja=standard,a4paper,11pt]{bxjsarticle}

\usepackage{amsmath,amssymb,amsthm,mathtools,bm}
\usepackage{booktabs,longtable,array}
\usepackage{geometry}
\usepackage{hyperref}
\usepackage{tikz}
\usetikzlibrary{arrows.meta,positioning}

\geometry{margin=25mm}
\hypersetup{
  unicode=true,
  colorlinks=true,
  linkcolor=blue,
  citecolor=blue,
  urlcolor=blue
}

\theoremstyle{definition}
\newtheorem{definition}{定義}[section]
\newtheorem{assumption}[definition]{仮定}
\newtheorem{remark}[definition]{注}
\theoremstyle{plain}
\newtheorem{lemma}[definition]{補題}
\newtheorem{theorem}[definition]{定理}
\newtheorem{proposition}[definition]{命題}
\newtheorem{corollary}[definition]{系}

\newcommand{\dd}{\,\mathrm{d}}
\newcommand{\E}{\mathbb{E}}
\newcommand{\Prob}{\mathbb{P}}
\newcommand{\Pareto}{\mathrm{Pareto}}
\newcommand{\Exp}{\mathrm{Exp}}
\newcommand{\Qfull}{Q_{\mathrm{full}}}
\newcommand{\Dstar}{\Delta_\infty^\ast}

\title{重尾潜伏下のランサムウェア攻撃に対する\\
攻撃・防疫合成半マルコフモデルと最適バックアップ設計\\
\large 検知投資の相転移と最適間隔の頻度漸近独立性}
\author{匿名査読用}
\date{}

\begin{document}
\maketitle

\begin{abstract}
本稿は，ランサムウェアの潜伏時間が
$T_d\sim\Pareto(t_m,\alpha)$ に従う環境で，攻撃フェーズ
（感染・潜伏・発症）と防疫フェーズ（検知・隔離・復旧）を
合成半マルコフ過程として定式化し，定常コストレートを最小化する
バックアップ設計問題を扱う．中心的な結論は四つである．
第一に，検知レート $\beta(I)>0$ が正ならば
$\E[\min(T_{\mathrm{det}},T_d)]$ はすべての $\alpha>0$ で有限となり，
更新報酬定理の適用を尾指数の制限なしに正当化できる．
第二に，$\alpha\in(0,1)$ での早期検知確率の Tauber 展開について，
剰余項が $O(\beta)$ であることを初等的不等式
$1-e^{-x}\le x$ から直接示せる．
第三に，完全モデルの期待損失 $\Qfull(\Delta)$ の唯一の大域最小点は
\[
  \Dstar=\frac{L-\eta(d_1\tau_R+d_0)}{\eta d_1\kappa_s n}
\]
であり，尾指数 $\alpha$，攻撃頻度 $\lambda$，検知レート $\beta(I)$
には依存しない．
第四に，固定費を含む最適間隔 $\Delta^\ast(\lambda)$ は
$\lambda\to\infty$ で $\Dstar$ に収束する．
さらに，動的保持延長を単期間費用最適化問題として定式化し，
閉形式の整数最適解を与える．
\end{abstract}

\tableofcontents

\section{序論}

ランサムウェア攻撃では，侵入から発症までの潜伏時間
（dwell time）が長く，かつ強い重尾性を持つことが報告されている．
指数分布の尾確率 $\Prob(T>t)=e^{-\mu t}$ は急速に減衰するが，
Pareto 分布
\[
  \Prob(T_d>t)=\left(\frac{t_m}{t}\right)^\alpha,\qquad t\ge t_m,
\]
はべき乗則で減衰する．特に $\alpha\le 1$ では
$\E[T_d]=\infty$ となり，有限平均寿命を暗黙に仮定する古典的な
予防保全モデルや更新報酬モデルをそのまま適用することが難しい．

本稿の問題意識は，重尾潜伏がバックアップ間隔と検知投資の最適化に
与える理論的影響を，半マルコフ過程と費用レート最小化の枠組みで
明確化することである．主な貢献は以下の通りである．

\begin{enumerate}
  \item 検知レートが正であれば，潜伏状態の平均滞在時間が
  すべての $\alpha>0$ で有限となることを示す．
  \item 早期検知確率の低投資域漸近で，$\alpha=1$ を境界とする
  相転移を導く．
  \item 完全モデルの期待損失 $\Qfull(\Delta)$ について，
  大域最小点 $\Dstar$ を閉形式で与える．
  \item 固定バックアップ費を含む最適間隔
  $\Delta^\ast(\lambda)$ が，高攻撃頻度極限で $\Dstar$ に
  収束することをコンパクト性から証明する．
\end{enumerate}

\section{モデルの定式化}

\subsection{半マルコフ過程}

半マルコフ過程は，遷移先は現在状態のみに依存するが，各状態での
滞在時間には指数分布に限らない任意の分布を許す確率過程である．
本稿では，潜伏時間に Pareto 分布を直接導入するため，この枠組みを
用いる．

\subsection{攻撃フェーズと防疫フェーズ}

\begin{definition}[攻撃フェーズ]
状態空間を $\{A_0,A_1,A_2,A_3\}$ とする．
$A_0$ は無感染，$A_1$ は潜伏感染，$A_2$ は発症中，
$A_3$ は破壊完了を表す．攻撃到着は強度 $\lambda>0$ の
Poisson 過程であり，攻撃到着時に $A_0\to A_1$ が起こる．
潜伏時間 $T_d$ の経過後に $A_1\to A_2$，さらに
$T_e\sim\Exp(\mu_e)$ の後に $A_2\to A_3$ が起こる．
\end{definition}

\begin{assumption}[Pareto 潜伏時間]
$T_d\sim\Pareto(t_m,\alpha(I))$ とし，
\[
  \Prob(T_d>t)=\left(\frac{t_m}{t}\right)^{\alpha(I)},\qquad
  f_{T_d}(t)=\alpha(I)t_m^{\alpha(I)}t^{-(\alpha(I)+1)}
\]
for $t\ge t_m>0$ とする．尾指数 $\alpha(I)>0$ は投資 $I$ の関数である．
\end{assumption}

\begin{definition}[防疫フェーズ]
状態空間を $\{D_0,D_1,D_2,D_3\}$ とする．
$D_0$ は通常，$D_1$ は検知，$D_2$ は隔離，$D_3$ は復旧中を表す．
早期検知時刻は $T_{\mathrm{det}}\sim\Exp(\beta(I))$ で，
$T_{\mathrm{det}}<T_d$ の場合に早期検知が生起する．発症後には
確率 $\eta\in(0,1]$ で症状検知される．
\end{definition}

\begin{assumption}[投資応答関数]
$\alpha(I)$ と $\beta(I)$ は $[0,\infty)$ 上の $C^2$ 関数で，
単調増加かつ凹型である．また
\[
  \lim_{I\to\infty}\alpha(I)=\alpha_{\max}<\infty,\qquad
  \lim_{I\to\infty}\beta(I)=\beta_{\max}<\infty,
\]
かつ $\alpha(0)>0,\ \beta(0)>0$ とする．
\end{assumption}

合成状態空間は
\[
S=\{S_N,S_L,S_E,S_S,S_I,S_R,S_F\}
\]
で表す．$S_N$ は正常，$S_L$ は潜伏・未検知，$S_E$ は早期検知済，
$S_S$ は発症後・症状検知待ち，$S_I$ は隔離済，$S_R$ は復旧中，
$S_F$ は全損を表す．

\begin{figure}[htbp]
\centering
\begin{tikzpicture}[
  node distance=16mm and 22mm,
  state/.style={draw,rounded corners,minimum width=13mm,minimum height=8mm,align=center},
  loss/.style={draw,rounded corners,fill=gray!20,minimum width=13mm,minimum height=8mm,align=center},
  arr/.style={-{Latex[length=2mm]},thick}
]
\node[state] (SN) {$S_N$};
\node[state,right=of SN] (SL) {$S_L$};
\node[state,above right=of SL] (SE) {$S_E$};
\node[state,right=of SE] (SI) {$S_I$};
\node[state,right=of SI] (SR) {$S_R$};
\node[state,below right=of SL] (SS) {$S_S$};
\node[loss,right=of SS] (SF) {$S_F$};
\draw[arr] (SN) -- node[above] {$\lambda$} (SL);
\draw[arr] (SL) -- node[left] {$\beta(I)$} (SE);
\draw[arr] (SE) -- node[above] {隔離} (SI);
\draw[arr] (SI) -- node[above] {復旧} (SR);
\draw[arr] (SL) -- node[left] {$T_d$ 後} (SS);
\draw[arr] (SS) -- node[above] {$1-\eta$} (SF);
\draw[arr] (SS) -- node[left] {$\eta$} (SI);
\draw[arr] (SR) to[bend left=35] node[above] {完了} (SN);
\end{tikzpicture}
\caption{合成半マルコフ状態遷移図．}
\end{figure}

\subsection{バックアップと復旧}

\begin{definition}[バックアップスケジュール]
間隔 $\Delta>0$ で $n\in\mathbb{Z}_+$ 世代のバックアップを保持する．
侵入後に取得されたバックアップは汚染されるとし，汚染世代数は
早期検知時に $k=\lceil T_{\mathrm{det}}/\Delta\rceil$，
症状検知時に $k=\lceil T_d/\Delta\rceil$ とする．
$k\ge n$ なら全損，$k<n$ なら復旧可能である．
\end{definition}

\begin{definition}[コンテナベース復旧時間]
復旧は隔離，クローン，再デプロイの三段階から成る．
$T_{\mathrm{iso}}\sim\Exp(\mu_{\mathrm{iso}})$，
$T_{\mathrm{clone}}\sim\Exp(\mu_c)$，
$T_{\mathrm{deploy}}\sim\Exp(\mu_d)$ とし，
\[
  \tau_R:=\mu_{\mathrm{iso}}^{-1}+\mu_c^{-1}+\mu_d^{-1}.
\]
全停止時間は $T_{\mathrm{down}}=T_{\mathrm{iso}}+T_{\mathrm{rec}}$，
$T_{\mathrm{rec}}=T_{\mathrm{scan}}+T_{\mathrm{clone}}+T_{\mathrm{deploy}}$，
$T_{\mathrm{scan}}=\kappa_sT_{\mathrm{det}}$ とする．
\end{definition}

\subsection{期待損失と費用レート}

更新サイクルあたりの期待損失を
\begin{equation}
  \Qfull(\Delta;\alpha)
  =(1-p_e)\left\{\eta(d_1\tau_R+d_0)(1-F)
  +\eta d_1\kappa_s ET+F L\right\}
  \label{eq:qfull}
\end{equation}
と定義する．ここで
\[
  p_e=\Prob(T_{\mathrm{det}}<T_d),\quad
  F(\Delta;\alpha)=\left(\frac{t_m}{n\Delta}\right)^\alpha,\quad
  ET(\Delta;\alpha)=\E[T_d\mathbf{1}_{\{T_d\le n\Delta\}}].
\]
バックアップ固定費，投資費，追加保持費を含む定常費用レートは
\begin{equation}
  C(\Delta,I,n_{\mathrm{ext}})
  =\lambda\Qfull(\Delta;\alpha(I))
  +\frac{c_b}{\Delta}+c_I I+c_n n_{\mathrm{ext}}.
  \label{eq:cost-rate}
\end{equation}

\section{更新報酬定理の適用条件}

\begin{lemma}[潜伏状態の平均滞在時間の有限性]
$t_m>0,\ \beta(I)>0,\ \alpha(I)>0$ の任意の組に対し，
\[
  \E[\min(T_{\mathrm{det}},T_d)]
  =\frac{1-e^{-\beta(I)t_m}}{\beta(I)}
  +t_m^{\alpha(I)}\beta(I)^{\alpha(I)-1}
  \Gamma(1-\alpha(I),\beta(I)t_m)
  <\infty .
  \label{eq:min-finite}
\]
\end{lemma}

\begin{proof}
$\E[\min(X,Y)]=\int_0^\infty\Prob(X>t,Y>t)\dd t$ と独立性を用いると，
\[
  \E[\min(T_{\mathrm{det}},T_d)]
  =\int_0^{t_m}e^{-\beta t}\dd t
  +\int_{t_m}^{\infty}e^{-\beta t}\left(\frac{t_m}{t}\right)^\alpha\dd t .
\]
第一項は $(1-e^{-\beta t_m})/\beta$ である．第二項で $u=\beta t$ と置換すると
$t_m^\alpha\beta^{\alpha-1}\Gamma(1-\alpha,\beta t_m)$ を得る．
$x_0=\beta t_m>0$ に対し，$u\ge x_0$ では $u^{-\alpha}\le x_0^{-\alpha}$ なので，
\[
  \Gamma(1-\alpha,x_0)
  =\int_{x_0}^{\infty}u^{-\alpha}e^{-u}\dd u
  \le x_0^{-\alpha}e^{-x_0}<\infty .
\]
従って右辺は有限である．
\end{proof}

\begin{theorem}[更新報酬定理の適用可能性]
仮定 2.2 と 2.4 のもとで，任意の $\lambda>0,\ \beta(I)>0,\ n\ge1,\ \Delta>0$
に対して，合成半マルコフ過程の非吸収部分は正再帰的であり，
サイクルあたり期待費用は有限である．従って，更新報酬定理を適用でき，
\eqref{eq:cost-rate} は長期平均費用レートと一致する．
\end{theorem}

\section{Tauber 漸近展開と相転移}

早期検知確率は
\[
  p_e(\beta)=\Prob(T_{\mathrm{det}}<T_d)=1-L_{T_d}(\beta)
\]
である．低投資域 $\beta\to0+$ での振る舞いは，潜伏時間分布の尾の重さを反映する．

\begin{theorem}[Tauber 展開の剰余項]
$T_d\sim\Pareto(t_m,\alpha)$，$\alpha\in(0,1)$，$\beta>0$ のとき，
\begin{equation}
  p_e(\beta)=t_m^\alpha\Gamma(1-\alpha)\beta^\alpha-C(\beta),
  \label{eq:tauber}
\end{equation}
ただし
\[
  C(\beta)=\alpha t_m^\alpha\int_0^{t_m}(1-e^{-\beta t})t^{-(\alpha+1)}\dd t
\]
であり，
\begin{equation}
  0\le C(\beta)\le \frac{\alpha t_m}{1-\alpha}\beta .
  \label{eq:tauber-bound}
\end{equation}
従って $R(\beta)=p_e(\beta)-t_m^\alpha\Gamma(1-\alpha)\beta^\alpha=O(\beta)$ で，
相対誤差は $O(\beta^{1-\alpha})\to0$ である．
\end{theorem}

\begin{proof}
独立性から
\[
  p_e(\beta)=\int_{t_m}^{\infty}(1-e^{-\beta t})
  \alpha t_m^\alpha t^{-(\alpha+1)}\dd t .
\]
積分区間 $[t_m,\infty)$ を $[0,\infty)-[0,t_m]$ に分解し，
$u=\beta t$ と置換すると主要項
$t_m^\alpha\Gamma(1-\alpha)\beta^\alpha$ を得る．
補正項については $1-e^{-x}\le x$ を用いて
\[
  C(\beta)\le \alpha t_m^\alpha\beta\int_0^{t_m}t^{-\alpha}\dd t
  =\frac{\alpha t_m}{1-\alpha}\beta .
\]
\end{proof}

\begin{lemma}[Pareto Laplace 変換の漸近]
$T_d\sim\Pareto(t_m,\alpha)$，$\beta\to0+$ において，
\[
\E[T_de^{-\beta T_d}]
\sim
\begin{cases}
\alpha t_m^\alpha\Gamma(1-\alpha)\beta^{\alpha-1}, & \alpha\in(0,1),\\
-t_m\log(\beta t_m), & \alpha=1,\\
\alpha t_m/(\alpha-1), & \alpha>1.
\end{cases}
\]
\end{lemma}

\begin{theorem}[純粋検知費用の相転移]
バックアップ世代境界なしの極限で
\[
  Q_\infty(\beta)=d_1\{\tau_RL_{T_d}(\beta)+\E[T_de^{-\beta T_d}]\}
\]
とおく．このとき $\beta\to0+$ で
\[
Q_\infty(\beta)\sim
\begin{cases}
d_1\alpha t_m^\alpha\Gamma(1-\alpha)\beta^{\alpha-1}\to\infty,
& \alpha\in(0,1),\\
-d_1t_m\log(\beta t_m)\to\infty,& \alpha=1,\\
d_1(\tau_R+\alpha t_m/(\alpha-1))<\infty,& \alpha>1.
\end{cases}
\]
\end{theorem}

\begin{remark}
$\alpha\le1$ では低投資域で純粋検知設計の期待費用が発散する．
従って，重尾潜伏下ではバックアップ設計が検知投資の補完として不可欠である．
\end{remark}

\section{最適バックアップ設計}

投資 $I$ と追加保持世代数を固定すると，$\Delta$ に関する最適化は
\[
  G(\Delta;\lambda)=\frac{c_b}{\Delta}+\lambda\Qfull(\Delta;\alpha)
\]
の最小化に帰着する．

\begin{assumption}
\[
  \Lambda:=L-\eta(d_1\tau_R+d_0)>0 .
\]
すなわち，全損時損失は期待復旧費用を上回る．
\end{assumption}

\begin{theorem}[$\Qfull(\Delta)$ の閉形式微分と唯一の大域最小点]
$\Qfull(\Delta;\alpha)$ を \eqref{eq:qfull} で定義する．
仮定 5.1 のもとで
\begin{equation}
  \frac{\partial\Qfull}{\partial\Delta}
  =(1-p_e)\alpha F(\Delta;\alpha)
  \left(\eta d_1\kappa_s n-\frac{\Lambda}{\Delta}\right)
  \label{eq:qfull-derivative}
\end{equation}
が成立する．従って $\partial\Qfull/\partial\Delta=0$ の唯一解は
\begin{equation}
  \Dstar=\frac{\Lambda}{\eta d_1\kappa_s n}
  =\frac{L-\eta(d_1\tau_R+d_0)}{\eta d_1\kappa_s n}.
  \label{eq:dstar}
\end{equation}
また $\Delta<\Dstar$ では微分は負，$\Delta>\Dstar$ では正であり，
$\Dstar$ は唯一の大域最小点である．さらに
\[
  \left.\frac{\partial^2\Qfull}{\partial\Delta^2}\right|_{\Dstar}
  =(1-p_e)\alpha\frac{F(\Dstar;\alpha)}{\Dstar}\eta d_1\kappa_s n>0 .
\]
\end{theorem}

\begin{proof}
$F(\Delta;\alpha)=t_m^\alpha n^{-\alpha}\Delta^{-\alpha}$ なので
$\partial F/\partial\Delta=-\alpha F/\Delta$．また Leibniz の積分則より
$\partial ET/\partial\Delta=\alpha nF$．
\eqref{eq:qfull} を微分すると
\[
  \frac{\partial\Qfull}{\partial\Delta}
  =(1-p_e)
  \left\{-\eta(d_1\tau_R+d_0)\frac{\partial F}{\partial\Delta}
  +\eta d_1\kappa_s\frac{\partial ET}{\partial\Delta}
  +L\frac{\partial F}{\partial\Delta}\right\},
\]
これを整理して \eqref{eq:qfull-derivative} を得る．
$(1-p_e)\alpha F>0$ なので符号は括弧内だけで決まり，
零点は \eqref{eq:dstar} の一点に限られる．
\end{proof}

\begin{remark}
$\Dstar$ は $\alpha,\lambda,\beta(I)$ に依存しない．ただし，
微分の大きさには $(1-p_e)\alpha F$ が含まれるため，収束速度や
有限 $\lambda$ での補正には尾指数や検知レートが影響する．
\end{remark}

\begin{theorem}[頻度漸近独立性]
$c_b>0$ とし，$G(\Delta;\lambda)=c_b/\Delta+\lambda\Qfull(\Delta;\alpha)$
の最小点を $\Delta^\ast(\lambda)$ と書く．このとき
\[
  \lim_{\lambda\to\infty}\Delta^\ast(\lambda)=\Dstar .
\]
\end{theorem}

\begin{proof}
最適性から
\[
  \frac{c_b}{\Delta^\ast}+\lambda\Qfull(\Delta^\ast)
  \le
  \frac{c_b}{\Dstar}+\lambda\Qfull(\Dstar)
\]
であり，従って
\[
  \Qfull(\Delta^\ast(\lambda))
  \le \Qfull(\Dstar)+\frac{c_b}{\lambda\Dstar}.
\]
任意の $\lambda_0>0$ に対し，
\[
  M_0:=\Qfull(\Dstar)+\frac{c_b}{\lambda_0\Dstar}
\]
とおけば，すべての $\lambda\ge\lambda_0$ で
$\Delta^\ast(\lambda)$ はレベル集合
$K_{M_0}=\{\Delta>0:\Qfull(\Delta)\le M_0\}$ に入る．
$\Delta\to0+$ では $F(\Delta;\alpha)\to\infty$ なので
$\Qfull(\Delta)\to\infty$．$\alpha\le1$ では
$ET(\Delta;\alpha)\to\infty$ により上側も有界となる．
$\alpha>1$ では $\Qfull(\Delta)$ の無限遠極限が
$\Qfull(\Dstar)$ より大きいことから，$\lambda_0$ を十分大きく取れば
同様に上側有界となる．従って $K_{M_0}$ はコンパクトである．

任意の部分列は収束部分列を持つ．一階条件
\[
  \lambda\frac{\partial\Qfull}{\partial\Delta}(\Delta^\ast(\lambda))
  =\frac{c_b}{(\Delta^\ast(\lambda))^2}
\]
を $\lambda$ で割り，コンパクト性による正の下界と
$\lambda\to\infty$ を用いると，極限点では
$\partial\Qfull/\partial\Delta=0$ が成立する．
定理 5.2 の一意性により極限点は必ず $\Dstar$ であるため，
全列が $\Dstar$ に収束する．
\end{proof}

\begin{corollary}[有限頻度での近似]
$Q''_\infty=\left.\partial^2\Qfull/\partial\Delta^2\right|_{\Dstar}$ とおくと，
一階条件の Taylor 展開から
\[
  \Delta^\ast(\lambda)
  =\Dstar+\frac{c_b}{\lambda{\Dstar}^2 Q''_\infty}+O(\lambda^{-2})
\]
を得る．
\end{corollary}

\section{動的保持延長}

\begin{definition}[追加保持コスト関数]
$n_{\mathrm{ext}}\in\mathbb{Z}_{\ge0}$，バックアップ間隔 $\Delta$，
実効保持期間
\[
  T(n_{\mathrm{ext}})=(n+n_{\mathrm{ext}})\Delta-\tau_R>t_m
\]
とする．追加保持コストを $c_n n_{\mathrm{ext}}$，
復旧不能ペナルティを
$L_{\mathrm{fail}}(t_m/T(n_{\mathrm{ext}}))^\alpha$ とすれば，
\[
  B(n_{\mathrm{ext}})
  =c_n n_{\mathrm{ext}}
  +L_{\mathrm{fail}}\left(\frac{t_m}{T(n_{\mathrm{ext}})}\right)^\alpha .
\]
\end{definition}

\begin{proposition}[最適追加保持世代数]
$B$ は $n_{\mathrm{ext}}\in\mathbb{R}_{\ge0}$ 上で狭義凸であり，
連続緩和の最適点は
\[
  T^\ast=\left(\frac{\alpha L_{\mathrm{fail}}t_m^\alpha\Delta}{c_n}
  \right)^{1/(\alpha+1)},\qquad
  n_{\mathrm{ext}}^c=\frac{T^\ast+\tau_R}{\Delta}-n .
\]
整数最適解は
\[
  n_{\mathrm{ext}}^\ast=\max(0,\lceil n_{\mathrm{ext}}^c\rceil)
\]
で与えられる．
\end{proposition}

\begin{proof}
$T=(n+n_{\mathrm{ext}})\Delta-\tau_R$ として微分すると
\[
  \frac{\dd B}{\dd n_{\mathrm{ext}}}
  =c_n-L_{\mathrm{fail}}\alpha t_m^\alpha T^{-(\alpha+1)}\Delta .
\]
一階条件から $T=T^\ast$ を得る．また
\[
  \frac{\dd^2 B}{\dd n_{\mathrm{ext}}^2}
  =L_{\mathrm{fail}}\alpha(\alpha+1)t_m^\alpha T^{-(\alpha+2)}\Delta^2>0
\]
なので狭義凸である．非負整数制約では連続最適点の上側近傍を取り，
負の場合は境界 $0$ が最適となる．
\end{proof}

\section{数値例}

社会インフラ・エンタープライズ規模を想定したベースラインとして，
$t_m=1$ 日，$\lambda=1.0$ 件/日，$d_1=1667$ 万円/日，
$L=50000$ 万円，$c_b=10$ 万円・日，$n=5$ 世代，
$\tau_R=2.0$ 日，$\eta=0.95$，$\kappa_s=0.3$ 日/世代を用いる．
このとき
\[
  \Lambda=50000-0.95(1666.7\times2.0)=46833,
  \qquad
  \Dstar=\frac{46833}{0.95\times1666.7\times0.3\times5}
  =19.72\ \text{日}.
\]

\begin{table}[htbp]
\centering
\caption{$\E[\min(T_{\mathrm{det}},T_d)]$（日，$t_m=1$）}
\begin{tabular}{cccccc}
\toprule
$\beta$ & $\alpha=0.3$ & $\alpha=0.5$ & $\alpha=0.8$ & $\alpha=1.0$ & $\alpha=1.5$\\
\midrule
0.1 & 6.086 & 4.621 & 3.309 & 2.775 & 2.027\\
0.5 & 1.720 & 1.582 & 1.426 & 1.347 & 1.205\\
1.0 & 0.942 & 0.911 & 0.873 & 0.852 & 0.810\\
\bottomrule
\end{tabular}
\end{table}

\begin{table}[htbp]
\centering
\caption{補正項 $C(\beta)$ と上界の比較（$t_m=1$）}
\begin{tabular}{ccccc}
\toprule
& \multicolumn{2}{c}{$\alpha=0.5$} & \multicolumn{2}{c}{$\alpha=0.8$}\\
$\beta$ & $C(\beta)$ & 上界 & $C(\beta)$ & 上界\\
\midrule
$10^{-3}$ & $9.998\times10^{-4}$ & $1.0\times10^{-3}$ & $3.9997\times10^{-3}$ & $4.0\times10^{-3}$\\
$10^{-2}$ & $9.983\times10^{-3}$ & $1.0\times10^{-2}$ & $3.997\times10^{-2}$ & $4.0\times10^{-2}$\\
$10^{-1}$ & $9.837\times10^{-2}$ & $1.0\times10^{-1}$ & $3.967\times10^{-1}$ & $4.0\times10^{-1}$\\
\bottomrule
\end{tabular}
\end{table}

\begin{table}[htbp]
\centering
\caption{$\Delta^\ast(\lambda)\to\Dstar=19.72$ 日の収束}
\begin{tabular}{ccccccc}
\toprule
& \multicolumn{5}{c}{$\lambda$} & \\
$\alpha$ & 0.01 & 0.1 & 1.0 & 10 & 100 & $\Dstar$\\
\midrule
0.5 & 20.19 & 19.77 & 19.72 & 19.720 & 19.719 & 19.72\\
0.7 & 20.55 & 19.80 & 19.73 & 19.720 & 19.719 & 19.72\\
1.0 & 22.06 & 19.95 & 19.74 & 19.722 & 19.720 & 19.72\\
1.5 & 42.44 & 21.33 & 19.88 & 19.735 & 19.721 & 19.72\\
\bottomrule
\end{tabular}
\end{table}

\begin{table}[htbp]
\centering
\caption{最適追加保持世代数（$\Delta=10$ 日，$c_n=50$ 万円/世代，$L_{\mathrm{fail}}=50000$ 万円）}
\begin{tabular}{ccccc}
\toprule
$\alpha$ & $T^\ast$（日） & $n_{\mathrm{ext}}^c$ & $n_{\mathrm{ext}}^\ast$ & $B(n_{\mathrm{ext}}^\ast)$\\
\midrule
0.5 & 292.4 & 24.44 & 25 & 4146.4\\
0.7 & 182.7 & 13.47 & 14 & 1979.6\\
1.0 & 100.0 & 5.20 & 6 & 763.0\\
1.5 & 46.8 & -0.12 & 0 & 150.4\\
\bottomrule
\end{tabular}
\end{table}

\section{まとめ}

本稿は，重尾潜伏時間を持つランサムウェア攻撃に対して，
攻撃・防疫合成半マルコフモデルを構成し，バックアップ間隔の
最適設計を解析した．検知レート $\beta(I)>0$ が正であれば，
$\alpha\le1$ でも潜伏状態の平均滞在時間は有限となり，更新報酬定理を
適用できる．また，Tauber 展開は $\alpha=1$ を境界とする相転移を示し，
重尾環境ではバックアップが検知投資の不可欠な補完になる．

完全モデルの閉形式微分から得られる
\[
  \Dstar=\frac{L-\eta(d_1\tau_R+d_0)}{\eta d_1\kappa_s n}
\]
は，尾指数，攻撃頻度，検知レートに依存しない．固定費を含む最適解も
攻撃頻度が高まるにつれてこの値へ収束するため，$\Dstar$ は
高頻度攻撃環境における設計基準として機能する．
今後の課題は，有限 $\lambda$ での収束速度の明示的上界，復旧時間分布の
位相型分布への拡張，および実データからの尾指数推定である．

\begin{thebibliography}{48}
\bibitem{mandiant2023} Mandiant (2023). M-Trends 2023: Special Report. Google Cloud.
\bibitem{ibm2023} IBM Security (2023). X-Force Threat Intelligence Index 2023. IBM Corporation.
\bibitem{verizon2023} Verizon (2023). Data Breach Investigations Report 2023. Verizon Communications.
\bibitem{sophos2023} Sophos (2023). The State of Ransomware 2023. Sophos Ltd.
\bibitem{crowdstrike2023} CrowdStrike (2023). 2023 Global Threat Report. CrowdStrike Inc.
\bibitem{ponemon2023} IBM Security / Ponemon Institute (2023). Cost of a Data Breach Report 2023. IBM Corporation.
\bibitem{gordonloeb2002} Gordon, L. A. and Loeb, M. P. (2002). The economics of information security investment. ACM Trans. Inf. Syst. Secur., 5(4), 438--457.
\bibitem{gordon2010} Gordon, L. A., Loeb, M. P., and Sohail, T. (2010). Market value of voluntary disclosures concerning information security. MIS Quarterly, 34(3), 567--594.
\bibitem{hausken2006} Hausken, K. (2006). Returns to information security investment. Inf. Syst. Frontiers, 8(5), 338--349.
\bibitem{laszka2017} Laszka, A., Farhang, S., and Grossklags, J. (2017). On the economics of ransomware. In Decision and Game Theory for Security, LNCS 10575, pp. 397--417. Springer.
\bibitem{fielder2016} Fielder, A., Panaousis, E., Malacaria, P., Hankin, C., and Smeraldi, F. (2016). Decision support approaches for cyber security investment. Decis. Support Syst., 86, 13--23.
\bibitem{moore2009} Moore, T., Clayton, R., and Anderson, R. (2009). The economics of online crime. J. Econ. Perspect., 23(3), 3--20.
\bibitem{bohme2010} Böhme, R. and Schwartz, G. (2010). Modeling cyber-insurance: Towards a unifying framework. WEIS 2010.
\bibitem{anderson2020} Anderson, R. (2020). Security Engineering: A Guide to Building Dependable Distributed Systems, 3rd ed. Wiley.
\bibitem{feller1971} Feller, W. (1971). An Introduction to Probability Theory and Its Applications, Vol. II, 2nd ed. Wiley.
\bibitem{embrechts1997} Embrechts, P., Klüppelberg, C., and Mikosch, T. (1997). Modelling Extremal Events for Insurance and Finance. Springer.
\bibitem{resnick2007} Resnick, S. I. (2007). Heavy-Tail Phenomena. Springer.
\bibitem{bingham1989} Bingham, N. H., Goldie, C. M., and Teugels, J. L. (1989). Regular Variation. Cambridge University Press.
\bibitem{dehaan2006} de Haan, L. and Ferreira, A. (2006). Extreme Value Theory. Springer.
\bibitem{pickands1975} Pickands, J. (1975). Statistical inference using extreme order statistics. Ann. Statist., 3(1), 119--131.
\bibitem{hill1975} Hill, B. M. (1975). A simple general approach to inference about the tail of a distribution. Ann. Statist., 3(5), 1163--1174.
\bibitem{mitzenmacher2004} Mitzenmacher, M. (2004). A brief history of generative models for power law and lognormal distributions. Internet Math., 1(2), 226--251.
\bibitem{klein2003} Klein, J. P. and Moeschberger, M. L. (2003). Survival Analysis, 2nd ed. Springer.
\bibitem{limnios2001} Limnios, N. and Opris an, G. (2001). Semi-Markov Processes and Reliability. Birkhäuser.
\bibitem{cinlar1975} Çınlar, E. (1975). Introduction to Stochastic Processes. Prentice-Hall.
\bibitem{howard1971} Howard, R. A. (1971). Dynamic Probabilistic Systems, Vol. II. Wiley.
\bibitem{cox1962} Cox, D. R. (1962). Renewal Theory. Methuen.
\bibitem{karlin1975} Karlin, S. and Taylor, H. M. (1975). A First Course in Stochastic Processes, 2nd ed. Academic Press.
\bibitem{asmussen2003} Asmussen, S. (2003). Applied Probability and Queues, 2nd ed. Springer.
\bibitem{ross2014} Ross, S. M. (2014). Introduction to Probability Models, 11th ed. Academic Press.
\bibitem{barlow1965} Barlow, R. E. and Proschan, F. (1965). Mathematical Theory of Reliability. Wiley.
\bibitem{nakagawa2005} Nakagawa, T. (2005). Maintenance Theory of Reliability. Springer.
\bibitem{aven1999} Aven, T. and Jensen, U. (1999). Stochastic Models in Reliability. Springer.
\bibitem{porteus2002} Porteus, E. L. (2002). Foundations of Stochastic Inventory Theory. Stanford University Press.
\bibitem{zipkin2000} Zipkin, P. H. (2000). Foundations of Inventory Management. McGraw-Hill.
\bibitem{silver1998} Silver, E. A., Pyke, D. F., and Peterson, R. (1998). Inventory Management and Production Planning and Scheduling, 3rd ed. Wiley.
\bibitem{rockafellar1970} Rockafellar, R. T. (1970). Convex Analysis. Princeton University Press.
\bibitem{boyd2004} Boyd, S. and Vandenberghe, L. (2004). Convex Optimization. Cambridge University Press.
\bibitem{rudin1976} Rudin, W. (1976). Principles of Mathematical Analysis, 3rd ed. McGraw-Hill.
\bibitem{abramowitz1965} Abramowitz, M. and Stegun, I. A. (1965). Handbook of Mathematical Functions. Dover.
\bibitem{olver2010} Olver, F. W. J. et al. (eds.) (2010). NIST Digital Library of Mathematical Functions. Cambridge University Press.
\bibitem{iso2019} ISO (2019). ISO 22301: Security and Resilience--Business Continuity Management Systems.
\bibitem{nist2010} NIST (2010). NIST SP 800-34 Rev. 1: Contingency Planning Guide for Federal Information Systems.
\bibitem{cisa2020} CISA (2020). Ransomware Guide. Cybersecurity and Infrastructure Security Agency.
\bibitem{nisc2023} 内閣サイバーセキュリティセンター (2023). 重要インフラのサイバーセキュリティに係る行動計画.
\bibitem{ipa2024} IPA (2024). 情報セキュリティ10大脅威 2024.
\bibitem{jpcert2023} JPCERT/CC (2023). インシデントハンドリング報告（2023年1月--12月）.
\bibitem{npa2023} 警察庁 (2023). 令和5年におけるサイバー空間をめぐる脅威の情勢等について.
\end{thebibliography}

\end{document}
